% Source: https://artofproblemsolving.com/community/c3h1082145_inequalities
% Stable identity: aops-topic-1082145-q01
\textbf{(a) Yes.} Extend $AP$ to meet $BC$ at $Q$.

\begin{asy}
usepackage("amsmath");
size(4cm);
defaultpen(fontsize(10pt));
pair A=(0,0),B=(5,0),C=(1.6,3.2),P=(2.1,1.1); pair Q=extension(A,P,B,C);
draw(A--B--C--cycle,linewidth(1)); draw(A--Q,dashed); draw(P--B,blue); draw(A--P,blue);
dot(P); dot(Q); label("$A$",A,SW);label("$B$",B,SE);label("$C$",C,N);label("$P$",P,SE);label("$Q$",Q,NE);
label("\mbox{}",(min(currentpicture,user=true)+max(currentpicture,user=true))/2,black); // keeps the text after the figure black when LaTeX typesets the labels inline
\end{asy}

By the triangle inequality, $PB<PQ+QB$, so $PA+PB<PA+PQ+QB=AQ+QB$. Also $AQ<AC+CQ$, so $AQ+QB<AC+CQ+QB=AC+CB$. Hence $PA+PB<CA+CB$.

\textbf{(b) No.} Let $ABC$ be equilateral with side $1$ and let $D$ be a point at a small height $h$ above its centre $O$. Then $DA+DB+DC=3\sqrt{\dfrac{1}{3}+h^2}\to\sqrt3\approx1.732$ as $h\to0$. Choose $h$ so small that $DA+DB+DC<1.8$. A point $P$ inside the tetrahedron very close to $A$ has $PA+PB+PC$ close to $0+1+1=2>1.8$. So the inequality fails for this tetrahedron.
